\section*{Bab \ref{ch:g8:circuit-laws} --- Hukum Arus dan Tegangan di dalam Rangkaian}

\begin{solution}{exo:g8:circuit-laws:1}
Hukum arus serinya: satu lingkar, satu \omterm{def:g8:intensity-ammeter:intensity}{kuat arus} di mana-mana. Hukum
percabangannya: arus masuk sama dengan arus keluar. Hukum \omterm{def:g8:voltage-voltmeter:voltage}{tegangan}
serinya: \omterm{def:g8:voltage-voltmeter:voltage}{tegangan} \omterm{def:g3:first-electric-circuit:battery}{baterainya} sama dengan jumlah bagian alatnya
mengelilingi lingkarnya. Hukum \omterm{def:g8:voltage-voltmeter:voltage}{tegangan} paralelnya: \omterm{def:g5:series-parallel-circuits:parallel}{cabang} yang
membentang pada percabangan yang sama berbagi satu \omterm{def:g8:voltage-voltmeter:voltage}{tegangan} yang sama.
\end{solution}

\begin{solution}{exo:g8:circuit-laws:2}
$0.8 - 0.55 = \qty{0.25}{A}$.
\end{solution}

\begin{solution}{exo:g8:circuit-laws:3}
$9 - 2.5 - 3.8 = \qty{2.7}{V}$.
\end{solution}

\begin{solution}{exo:g8:circuit-laws:4}
\qty{12}{V} masing-masing --- menurut hukum \omterm{def:g8:voltage-voltmeter:voltage}{tegangan} paralelnya. Hukum
yang sama menaruh \qty{230}{V} yang serupa pada setiap stopkontak
pemasangan kabel paralel rumahnya.
\end{solution}

\begin{solution}{exo:g8:circuit-laws:5}
Misalnya $I_C$: begitu $I_{\text{baterai}}$ dan $I_B$ diketahui,
hukum percabangannya memancangkan $I_C = \qty{0.2}{A}$ --- alat mana
pun yang dikirim ke sana hanya dapat menegaskannya (demikian pula
$U_C$, yang dipancangkan hukum \omterm{def:g8:voltage-voltmeter:voltage}{tegangan} paralelnya begitu $U_B$
diketahui).
\end{solution}

\begin{solution}{exo:g8:circuit-laws:6}
Lampu \qty{2.7}{V}-nya: pada arus yang sama, bagian \omterm{def:g8:voltage-voltmeter:voltage}{tegangan} yang
lebih besar menandai penentang yang lebih kuat. Bab berikutnya
\omterm{def:g6:measurement-in-science:measuring}{mengukur} perlawanan itu --- yaitu hambatan.
\end{solution}

\begin{solution}{exo:g8:circuit-laws:7}
$2.2 + 1.9 + 1.4 = \qty{5.5}{V}$ terhadap dompet \qty{6}{V}: setengah
\omterm{def:g8:voltage-voltmeter:voltage}{volt} melenceng --- di luar serakan yang jujur sebuah alat ukur untuk
bacaan seperti itu; satu papan ukurnya salah dibaca (atau ada satu
bagian yang tidak terukur).
\end{solution}

\begin{solution}{exo:g8:circuit-laws:8}
$0.4 + 0.3 + 0.3 = \qty{1.0}{A}$ yang pergi terhadap \qty{0.9}{A}
yang tiba: \omterm{def:g5:series-parallel-circuits:parallel}{cabang} jalannya ``menciptakan'' \qty{0.1}{A} --- dihukum
oleh hukum percabangannya.
\end{solution}

\begin{solution}{exo:g8:circuit-laws:9}
Belanja lingkarnya harus berjumlah \qty{4.5}{V} milik \omterm{def:g3:first-electric-circuit:battery}{baterainya}:
\omterm{ex:g3:first-electric-circuit:switch}{sakelar} terbuka \qty{4.5}{V} $+$ lampu menganggur \qty{0}{V} $= 4.5$
--- seimbang. (Sesudah \omterm{def:g3:sound-around-us:sound}{bunyi} kliknya bagiannya bertukar: \omterm{ex:g3:first-electric-circuit:switch}{sakelarnya}
dekat nol, lampunya dekat \qty{4.5}{V} --- seimbang lagi.)
\end{solution}

\begin{solution}{exo:g8:circuit-laws:10}
Sebelumnya: $3 \times 1.5 = \qty{4.5}{V}$. Sesudahnya: $4.5 - 1.5 =
\qty{3.0}{V}$ --- sel yang terbalik mengurangi. Lampunya meredup dari
nyala \qty{4.5}{V}-nya menjadi nyala \qty{3}{V}.
\end{solution}

\begin{solution}{exo:g8:circuit-laws:11}
\omterm{def:g5:series-parallel-circuits:parallel}{Cabang} satu membentang pada \omterm{def:g3:first-electric-circuit:battery}{baterainya}: $U_E + U_F = \qty{6}{V}$
(hukum \omterm{def:g8:voltage-voltmeter:voltage}{tegangan} seri di sepanjang \omterm{def:g5:series-parallel-circuits:parallel}{cabangnya}), sehingga $U_F = 6 - 2.2 =
\qty{3.8}{V}$. \omterm{def:g5:series-parallel-circuits:parallel}{Cabang} dua membentang pada percabangan yang sama: $U_G =
\qty{6}{V}$ (hukum \omterm{def:g8:voltage-voltmeter:voltage}{tegangan} paralelnya).
\end{solution}

\begin{solution}{exo:g8:circuit-laws:12}
$I_R = 0.6 - 0.45 = \qty{0.15}{A}$ (hukum percabangannya). $U_R = U_Q
= \qty{2.3}{V}$ (hukum \omterm{def:g8:voltage-voltmeter:voltage}{tegangan} paralelnya). $U_P = 6 - 2.3 =
\qty{3.7}{V}$ (hukum \omterm{def:g8:voltage-voltmeter:voltage}{tegangan} serinya). Sebuah tambahan yang dapat
dihukum: ``$I_R = \qty{0.2}{A}$'' --- pembukuan percabangannya akan
menunjukkan $0.45 + 0.2 = 0.65 \neq 0.6$; atau ``$U_R = \qty{2.6}{V}$''
--- \omterm{def:g5:series-parallel-circuits:parallel}{cabang} paralel tidak boleh berselisih.
\end{solution}

\begin{solution}{pb:g8:circuit-laws:1}
\textbf{1.} Seri. Ketika \omterm{def:g4:series-circuits:series}{terangkai seri} kedua lampunya melawan
arak-arakannya satu demi satu: arus lingkarnya jatuh jauh di bawah
\qty{0.4}{A} milik satu lampu --- \qty{0.2}{A} yang terukur itu cocok.
(Paralel justru akan \emph{menjumlahkan} dua \omterm{def:g5:series-parallel-circuits:parallel}{cabang} penuh.)
\textbf{2.} Paralel: dua \omterm{def:g5:series-parallel-circuits:parallel}{cabang} yang masing-masing kira-kira
\qty{0.4}{A} --- kasarnya \qty{0.8}{A} dari \omterm{def:g3:first-electric-circuit:battery}{baterainya}.
\textbf{3.} Hukum \omterm{def:g8:voltage-voltmeter:voltage}{tegangan} serinya membelah \qty{6}{V}-nya di antara
dua lampu yang serupa: kira-kira \qty{3}{V} masing-masing.
\textbf{4.} Paralel --- tiap lampu akan membentang pada \qty{6}{V}
yang penuh. Lewat lubang intipnya pasangan serinya menyala redup:
masing-masing kurang makan pada separuh dorongan biasanya.
\textbf{5.} $I_Z = 0.5 - 0.35 = \qty{0.15}{A}$ (hukum
percabangannya); $I_X = \qty{0.5}{A}$ --- X duduk pada jalan yang tidak
terbagi (hukum arus serinya).
\textbf{6.} $U_{\text{pasangan}} = 6 - 2.8 = \qty{3.2}{V}$ (hukum
lingkarnya); $U_Y = U_Z = \qty{3.2}{V}$ (hukum paralelnya).
\textbf{7.} Dorongan yang sama, arak-arakan yang berbeda: kedua
lampunya melawan arusnya secara tidak sama --- Z, yang membawa lebih
sedikit pada \omterm{def:g8:voltage-voltmeter:voltage}{tegangan} yang sama, adalah penentang yang lebih kuat
(jalan yang lebih berat).
\textbf{8.} Turun: kehilangan \omterm{def:g5:series-parallel-circuits:parallel}{cabang} Y menyingkirkan sebuah jalan;
jalan Z yang tersisa melawan lebih daripada pasangannya bersama-sama,
sehingga \omterm{def:g3:first-electric-circuit:battery}{baterainya} menggerakkan arus total yang lebih sedikit.
\textbf{9.} Hukuman satu, hukum lingkarnya: $2.5 + 4.1 =
\qty{6.6}{V}$ dibelanjakan dari dompet \qty{6}{V}. Hukuman dua,
hukum percabangannya: masuk \qty{0.3}{A}, keluar $0.2 + 0.15 =
\qty{0.35}{A}$.
\textbf{10.} $6 - 2.5 = \qty{3.5}{V}$.
\textbf{11.} $0.2 + 0.15 = \qty{0.35}{A}$.
\textbf{12.} Karena keempat hukumnya menentukan berlebih setiap
rangkaian yang jujur: tiap bacaan harus sepakat dengan jumlah bacaan
yang lain, sehingga satu bilangan yang dikarang pasti berdusta kepada
sedikitnya satu hukum --- lalu dihukum oleh hitungan belaka.
\textbf{13.} Misalnya: ``Satu jalan, satu arak-arakan --- itulah hukum
arus serinya. Pada tiap \omterm{def:g5:series-parallel-circuits:parallel}{cabang} jalan arak-arakannya seimbang, masuk
sama dengan keluar --- itulah hukum percabangannya. Mengelilingi
lingkar mana pun, air terjun \omterm{def:g3:first-electric-circuit:battery}{baterainya} dibelanjakan sampai tetes
terakhir oleh pekerjanya --- itulah hukum \omterm{def:g8:voltage-voltmeter:voltage}{tegangan} serinya. Dan \omterm{def:g5:series-parallel-circuits:parallel}{cabang}
yang tergantung di antara dua percabangan yang sama minum dari tinggi
yang sama --- itulah hukum \omterm{def:g8:voltage-voltmeter:voltage}{tegangan} paralelnya.''
\end{solution}
